\chapter*{ВВЕДЕНИЕ}
\addcontentsline{toc}{chapter}{ВВЕДЕНИЕ}

\todo{Актуальность темы: связать управление спектрами низкоразмерных наноструктур с размерным квантованием и геометрией квантовых точек.}
\todo{Связь с физикой наноматериалов и нанотехнологий: подчеркнуть модельный характер работы и роль вычислительного нанодизайна.}
\todo{Объект исследования: модельные квантовые точки на квадратной tight-binding решётке.}
\todo{Предмет исследования: низкоэнергетические уровни и простые суррогатные аппроксимации для суперэллиптических геометрий.}

\textbf{Черновая формулировка цели.}
Целью работы является исследование возможности построения физически мотивированной суррогатной модели для аппроксимации низкоэнергетического спектра модельных квантовых точек с суперэллиптической геометрией на квадратной tight-binding решётке.

\textbf{Черновой список задач.}
\begin{enumerate}
  \item Implement and validate the tight-binding spectral calculation for finite model quantum dots.
  \item Construct fixed discrete-n superellipse datasets.
  \item Verify the physical energy convention and low-energy scaling.
  \item Build and evaluate a physics-informed Ridge surrogate.
  \item Compare Ridge with a tiny MLP ablation under structured validation.
  \item Analyze Ridge residuals against edge-discretization diagnostics.
  \item Optionally demonstrate surrogate-assisted inverse screening with direct Kwant validation.
\end{enumerate}

\todo{Методы исследования: tight-binding Hamiltonian, Kwant, структурированная кросс-валидация, Ridge-регрессия, диагностические корреляции.}
\todo{Научная новизна: сформулировать осторожно как контролируемый workflow для модельных суперэллиптических квантовых точек, не как новую физическую модель решёточных бильярдов.}
\todo{Практическая значимость: быстрый интерпретируемый суррогат внутри проверенного окна параметров; не утверждать замену Kwant.}
\todo{Достоверность результатов: прямой расчёт Kwant, физические sanity checks, предрегистрация критериев, структурированная валидация LOAO/LOARO.}
\todo{Структура работы: кратко описать главы после финализации текста.}

\clearpage
